
% !TEX TS-program = xelatex
\documentclass[12pt,a4paper]{article}
\usepackage{fontspec}
\usepackage{polyglossia}
\setmainlanguage{russian}
\setotherlanguage{english}
\setmainfont{DejaVu Serif}
\setsansfont{DejaVu Sans}
\setmonofont{DejaVu Sans Mono}
\usepackage[a4paper,margin=2.05cm]{geometry}
\usepackage{amsmath,amssymb,mathtools}
\usepackage{graphicx}
\usepackage{booktabs}
\usepackage{longtable}
\usepackage{array}
\usepackage{hyperref}
\usepackage{caption}
\usepackage{float}
\usepackage{enumitem}
\usepackage{xcolor}
\usepackage{microtype}
\hypersetup{colorlinks=true,linkcolor=blue,urlcolor=blue,citecolor=blue}
\setlength{\parskip}{0.55em}
\setlength{\parindent}{0pt}
\renewcommand{\arraystretch}{1.15}
\newcommand{\TMI}{\mathrm{TMI}}
\newcommand{\std}{\mathrm{std}}
\newcommand{\base}{\mathrm{base}}
\newcommand{\Prob}{\operatorname{Pr}}
\newcommand{\Auto}{\operatorname{Auto}}
\newcommand{\Safe}{\operatorname{Safe}}

\title{\textbf{Computational Validation of Adaptive Interface Control in TMI Quantum Systems}\\
\large Расширенный технический отчёт по вычислительным тестам 1--11}
\author{Salkutsan Aleksey}
\date{6 мая 2026}

\begin{document}
\maketitle

\begin{abstract}
В этом документе объединена вычислительная линия проверки предсказаний Теории многообразных интерфейсов (ТМИ) для квантовых систем. Версия отчёта расширяет первоначальные Tests 1--6 и добавляет Tests 7--11: полную density-matrix Monte Carlo проверку, cross-model validation, карту границ отказа, adversarial failure test и safety-aware recovery test.

Проверяемое инженерное ядро имеет форму
\[
I_Q^{coh-auto}\Rightarrow Loss_{coh}\downarrow\Rightarrow T_2^{eff}\uparrow\Rightarrow P_{success}\uparrow.
\]
Отчёт не утверждает лабораторное доказательство существования физического поля $\Xi_I^{auto}$. Он фиксирует более скромный результат: в вычислительных моделях адаптивное интерфейсное управление может увеличивать эффективную когерентность и вероятность правильного результата, но имеет границы применимости и требует safety-aware слоя в adversarial условиях.
\end{abstract}

\tableofcontents
\newpage

\section{Статус и научная осторожность}

Данный документ является \textbf{computational validation report}, а не лабораторным подтверждением. Мы проверяем не всю ТМИ сразу, а узкое следствие:
\[
\boxed{
I_Q^{coh-auto}
\Rightarrow
\gamma_{eff}\downarrow
\Rightarrow
T_2^{eff}\uparrow
\Rightarrow
P_{success}\uparrow
}
\]
где $I_Q^{coh-auto}$ -- когерентностно-сохраняющий автоматический интерфейс, $\gamma_{eff}=1/T_2^{eff}$ -- эффективная скорость декогеренции, а $P_{success}$ -- вероятность правильного результата малой квантовой схемы.

Начиная с Test 10 и Test 11, отчёт также фиксирует отрицательное и защитное следствие:
\[
\boxed{\Auto(I)\not\Rightarrow \Safe(I)}
\]
и
\[
\boxed{\Auto(I)+\text{SafetyLayer}\Rightarrow \text{SaferAdaptiveControl}.}
\]
То есть адаптивность сама по себе не является гарантией безопасности: при ложной обратной связи и высокой цене управления автоинтерфейс может проигрывать фиксированному стандартному режиму.

\section{Связь с обновлённым ядром ТМИ и TMI-Physics 2.2}

Тестовая линия не заменяет Core 0.2 и TMI-Physics 2.2. Её правильное место -- прикладной вычислительный модуль:
\[
\boxed{
\text{TMI Core 0.2}
\Rightarrow
\text{TMI-Physics 2.2}
\Rightarrow
\text{Computational Validation}
}
\]

В Core 0.2 автоматичность понимается как дефиниционное свойство действительного интерфейса:
\[
\operatorname{Valid}(I)\Rightarrow \Auto(I),
\qquad
\Auto(I)\not\Rightarrow \operatorname{AlwaysAct}(I).
\]
В TMI-Physics 2.2 эта логика переносится на физический слой:
\[
\Xi_I\in\mathfrak I_{phys}
\Rightarrow
\Auto(\Xi_I)
\Rightarrow
\Xi_I^{auto}.
\]
Вычислительные тесты проверяют инженерную проекцию этого принципа: может ли автоинтерфейс управления когерентностью улучшать наблюдаемые метрики в моделях квантовых систем.

\section{Сводка Tests 1--11}

В последнем столбце для Tests 1--4 порядок значений обычно: base / standard / TMI.

\begin{longtable}{p{2.0cm}p{3.9cm}p{4.2cm}p{5.0cm}}
\toprule
Тест & Смысл & Критерий & Результат \\
\midrule
\endfirsthead
\toprule
Тест & Смысл & Критерий & Результат \\
\midrule
\endhead

Test 1 & Ramsey toy model & $T_2^{TMI}>T_2^{std}>T_2^{base}$ & 50.0 / 120.0 / 178.6 $\mu s$ \\
Test 2 & Noisy Ramsey learning & $T_2^{TMI}>T_2^{std}>T_2^{base}$ & 51.3 / 125.8 / 162.4 $\mu s$ \\
Test 3 & Small circuit & $P_{TMI}>P_{std}>P_{base}$ & 0.182 / 0.386 / 0.431 \\
Test 4 & Stabilized interface & $P_{stab}>P_{std}>P_{base}$ & 0.180 / 0.328 / 0.465 \\
Test 5 & Surrogate Monte Carlo & $E[P_{TMI}]>E[P_{std}]$ & mean $\Delta=0.0512$, $Pr=0.864$ \\
Test 6 & Ablation study & full TMI $>$ standard & full $\Delta\approx 0.0677$ \\
Test 7 & Full density-matrix MC & $E[P_{TMI}]>E[P_{std}]$ & mean $\Delta=0.0935$, $Pr=0.933$ \\
Test 8 & Cross-model validation & advantage across circuit families & mean $\Delta=0.0662$, $Pr=0.846$ \\
Test 9 & Failure-boundary map & locate weak/failure regions & failure cells $=0$, weak cells $=6$ \\
Test 10 & Adversarial failure & force $P_{TMI}\le P_{std}$ & failure cells $=58$, weak cells $=22$ \\
Test 11 & Safety recovery & reduce failure rate & $\Delta_{safety}=0.0139$, recovered cells $=31/58$ \\

\bottomrule
\end{longtable}

\section{Tests 1--6: базовая линия вычислительной валидации}

\subsection{Test 1: Ramsey toy model}

Модель когерентности:
\[
C(t)=C_0 e^{-t/T_2}.
\]
Результат:
\[
T_2^{base}=50.0\,\mu s,
\quad
T_2^{std}=120.0\,\mu s,
\quad
T_2^{TMI}=178.6\,\mu s.
\]

\begin{figure}[H]
\centering
\includegraphics[width=0.86\linewidth]{figures/fig_test1_ramsey.png}
\caption{Test 1: простая Ramsey-модель когерентности.}
\end{figure}

\subsection{Test 2: Noisy Ramsey learning}

Автоинтерфейс должен учить параметр защиты $u$ по шумным Ramsey-наблюдениям. Результат:
\[
T_2^{base}=51.3\,\mu s,
\quad
T_2^{std}=125.8\,\mu s,
\quad
T_2^{TMI}=162.4\,\mu s.
\]

\begin{figure}[H]
\centering
\includegraphics[width=0.86\linewidth]{figures/fig_test2_noisy_ramsey.png}
\caption{Test 2: Ramsey-кривые при шумных наблюдениях.}
\end{figure}

\begin{figure}[H]
\centering
\includegraphics[width=0.86\linewidth]{figures/fig_test2_learning.png}
\caption{Test 2: обучение параметра защиты $u$.}
\end{figure}

\subsection{Test 3: Small quantum circuit}

Схема строится как
\[
|000\rangle\to \text{entangling circuit}\to \text{inverse circuit}\to |000\rangle.
\]
Идеальная вероятность успеха:
\[
P_{success}=\Pr(000).
\]
Результат:
\[
P_{base}=0.182,\quad P_{std}=0.386,\quad P_{TMI}=0.431.
\]

\begin{figure}[H]
\centering
\includegraphics[width=0.82\linewidth]{figures/fig_test3_success.png}
\caption{Test 3: вероятность правильного результата малой квантовой схемы.}
\end{figure}

\subsection{Test 4: Stabilized adaptive interface}

В Test 4 добавлены память, сглаживание, ограничение шага и повторные наблюдения:
\[
\text{AutoInterface}+\text{Memory}+\text{Smoothing}+\text{StepLimit}\Rightarrow\text{StableControl}.
\]
Результат:
\[
P_{base}=0.180,\quad P_{std}=0.328,\quad P_{raw}=0.465,\quad P_{stab}=0.465.
\]

\begin{figure}[H]
\centering
\includegraphics[width=0.82\linewidth]{figures/fig_test4_success.png}
\caption{Test 4: stabilized adaptive interface улучшает $P_{success}$.}
\end{figure}

\begin{figure}[H]
\centering
\includegraphics[width=0.86\linewidth]{figures/fig_test4_learning.png}
\caption{Test 4: raw adaptive скачет, stabilized adaptive удерживает параметр ближе к оптимальному режиму.}
\end{figure}

\subsection{Test 5: Monte Carlo surrogate validation}

Test 5 проверяет статистическое преимущество на 500 случайных средах с fast circuit-level surrogate:
\[
\mathbb E[P_{TMI}]>\mathbb E[P_{std}],
\qquad
\Pr(P_{TMI}>P_{std})>0.5.
\]
Результат:
\[
\mathbb E[P_{base}]=0.245,\quad
\mathbb E[P_{std}]=0.417,\quad
\mathbb E[P_{TMI}]=0.468.
\]
\[
\mathbb E[P_{TMI}-P_{std}]=0.0512,
\quad
\Pr(P_{TMI}>P_{std})=0.864.
\]

\begin{figure}[H]
\centering
\includegraphics[width=0.82\linewidth]{figures/fig_test5_mean.png}
\caption{Test 5: средний $P_{success}$ по Monte Carlo surrogate.}
\end{figure}

\begin{figure}[H]
\centering
\includegraphics[width=0.82\linewidth]{figures/fig_test5_delta.png}
\caption{Test 5: распределение выигрыша $P_{TMI}-P_{std}$.}
\end{figure}

\subsection{Test 6: Ablation study}

Test 6 показывает, что простая сырая адаптивность не гарантирует выигрыш:
\[
\text{RawAdaptivity}\not\Rightarrow\text{Improvement}.
\]
Выигрыш появляется как свойство структурированного автоинтерфейса:
\[
\text{StructuredAutoInterface}\Rightarrow\text{StableParameterTracking}\Rightarrow P_{success}\uparrow.
\]

\begin{figure}[H]
\centering
\includegraphics[width=0.90\linewidth]{figures/fig_test6_ablation.png}
\caption{Test 6: ablation study по конфигурациям интерфейсного управления.}
\end{figure}

\begin{figure}[H]
\centering
\includegraphics[width=0.90\linewidth]{figures/fig_test6_delta_full.png}
\caption{Test 6: потери относительно полного TMI-механизма.}
\end{figure}

\section{Tests 7--8: robustness и cross-model validation}

\subsection{Test 7: Full density-matrix Monte Carlo}

Test 7 повторяет статистическую проверку Test 5, но уже на полной 3-кубитной density-matrix toy-модели. Это устраняет главный недостаток Test 5: зависимость от fast surrogate.

\[
\mathbb E[P_{base}]=0.182,\quad
\mathbb E[P_{std}]=0.322,\quad
\mathbb E[P_{TMI}]=0.415.
\]
\[
\mathbb E[P_{TMI}-P_{std}]=0.0935,
\quad
\Pr(P_{TMI}>P_{std})=0.933.
\]

\begin{figure}[H]
\centering
\includegraphics[width=0.82\linewidth]{figures/fig_test7_mean.png}
\caption{Test 7: full density-matrix Monte Carlo, средний $P_{success}$.}
\end{figure}

\begin{figure}[H]
\centering
\includegraphics[width=0.70\linewidth]{figures/fig_test7_paired.png}
\caption{Test 7: paired comparison $P_{std}$ vs $P_{TMI}$.}
\end{figure}

\subsection{Test 8: Cross-model validation}

Test 8 проверяет не одну схему, а четыре circuit families. Общий результат:
\[
\mathbb E[P_{base}]=0.292,\quad
\mathbb E[P_{std}]=0.379,\quad
\mathbb E[P_{TMI}]=0.445.
\]
\[
\mathbb E[P_{TMI}-P_{std}]=0.0662,
\quad
\Pr(P_{TMI}>P_{std})=0.846.
\]

\begin{longtable}{p{5.3cm}rrrrr}
\toprule
Circuit & Base & Standard & TMI & $\Delta$ & $Pr$ \\
\midrule
\endfirsthead
\toprule
Circuit & Base & Standard & TMI & $\Delta$ & $Pr$ \\
\midrule
\endhead
balanced\_entangle\_inverse & 0.169 & 0.277 & 0.376 & 0.0994 & 1.000 \\
deep\_chain\_inverse & 0.253 & 0.322 & 0.377 & 0.0553 & 0.923 \\
ghz\_long\_idle\_inverse & 0.495 & 0.512 & 0.527 & 0.0148 & 0.462 \\
local\_mix\_inverse & 0.249 & 0.404 & 0.500 & 0.0955 & 1.000 \\
\bottomrule
\end{longtable}

\begin{figure}[H]
\centering
\includegraphics[width=0.95\linewidth]{figures/fig_test8_mean_by_circuit.png}
\caption{Test 8: средний $P_{success}$ по схемам.}
\end{figure}

\begin{figure}[H]
\centering
\includegraphics[width=0.90\linewidth]{figures/fig_test8_delta_by_circuit.png}
\caption{Test 8: выигрыш TMI над standard по схемам.}
\end{figure}

\paragraph{Важный нюанс.}
Для схемы \texttt{ghz\_long\_idle\_inverse} выигрыш в среднем положительный, но слабее и менее стабилен, чем на остальных схемах. Это показывает, что механизм не одинаково силён для всех типов схем.

\section{Tests 9--11: границы применимости и safety recovery}

\subsection{Test 9: Failure-boundary test}

Test 9 ищет не только успех, но и зоны ослабления. Мы варьируем скорость дрейфа скрытого оптимума и шум измерения. Критерий отказа:
\[
\mathbb E[P_{TMI}-P_{std}]\le 0
\quad\text{или}\quad
\Pr(P_{TMI}>P_{std})\le 0.5.
\]

Общий результат:
\[
\mathbb E[P_{TMI}-P_{std}]=0.0606,
\quad
\Pr(P_{TMI}>P_{std})=0.849.
\]
Найдено строгих failure cells:
\[
0,
\]
слабых зон:
\[
6.
\]

\begin{figure}[H]
\centering
\includegraphics[width=0.82\linewidth]{figures/fig_test9_delta_heatmap.png}
\caption{Test 9: heatmap выигрыша $P_{TMI}-P_{std}$ по drift/noise.}
\end{figure}

\begin{figure}[H]
\centering
\includegraphics[width=0.82\linewidth]{figures/fig_test9_prob_heatmap.png}
\caption{Test 9: вероятность $\Pr(P_{TMI}>P_{std})$ по drift/noise.}
\end{figure}

\subsection{Test 10: Adversarial failure test}

Test 10 специально создаёт враждебные условия:
\[
\text{fast drift}+\text{high noise}+\text{misleading feedback}+\text{high control penalty}.
\]
Цель -- добиться реальных случаев:
\[
P_{TMI}\le P_{std}.
\]

Результат:
\[
\mathbb E[P_{TMI}-P_{std}]=0.0025,
\quad
\Pr(P_{TMI}>P_{std})=0.530.
\]
Failure cells:
\[
58,
\]
weak cells:
\[
22.
\]

Ключевой вывод:
\[
\boxed{\text{AdaptiveInterface}\not\Rightarrow\text{SafeInterface}.}
\]

\begin{figure}[H]
\centering
\includegraphics[width=0.84\linewidth]{figures/fig_test10_failure_dist.png}
\caption{Test 10: adversarial failure distribution.}
\end{figure}

\begin{figure}[H]
\centering
\includegraphics[width=0.70\linewidth]{figures/fig_test10_failure_heatmap.png}
\caption{Test 10: число failure cells по misleading bias и control penalty.}
\end{figure}

\subsection{Test 11: Safety-aware recovery}

Test 11 берёт только failure cells из Test 10 и проверяет, можно ли восстановить преимущество через защитный слой:
\[
\text{AutoInterface}+	ext{SafetyPenalty}+	ext{ConservativeFallback}+	ext{ValidationChannel}\Rightarrow\text{SaferControl}.
\]

\begin{longtable}{p{3.8cm}rrrrr}
\toprule
Mode & Mean $P$ & $\Delta$ vs std & $Pr(>std)$ & Failure rate & Mean $u$ \\
\midrule
\endfirsthead
\toprule
Mode & Mean $P$ & $\Delta$ vs std & $Pr(>std)$ & Failure rate & Mean $u$ \\
\midrule
\endhead
standard & 0.2799 & 0.00000 & 0.000 & -- & 0.450 \\
adversarial\_TMI & 0.2799 & -0.00002 & 0.478 & 0.522 & 0.472 \\
safety\_v1 & 0.2939 & 0.01392 & 0.617 & 0.383 & 0.395 \\
recovery\_v2 & 0.2932 & 0.01328 & 0.579 & 0.421 & 0.386 \\
\bottomrule
\end{longtable}

Основной результат:
\[
\Delta_{adv}=-0.00002,
\quad
\Delta_{safety}=0.01392,
\quad
\Delta_{recovery}=0.01328.
\]
Recovery v2 восстановил:
\[
31/58
\]
failure cells.

\begin{figure}[H]
\centering
\includegraphics[width=0.82\linewidth]{figures/fig_test11_boxplot.png}
\caption{Test 11: safety recovery на зонах отказа из Test 10.}
\end{figure}

\begin{figure}[H]
\centering
\includegraphics[width=0.82\linewidth]{figures/fig_test11_failure_rate.png}
\caption{Test 11: снижение failure rate при safety-aware управлении.}
\end{figure}

\section{Итоговые теоретические следствия}

\subsection{Положительное следствие}

В обычных и умеренно шумных условиях вычислительная линия поддерживает инженерную гипотезу:
\[
\boxed{
I_Q^{coh-auto}
\Rightarrow
T_2^{eff}\uparrow
\Rightarrow
P_{success}\uparrow
}
\]

\subsection{Отрицательное следствие}

Test 10 вводит важную границу:
\[
\boxed{
\text{MisleadingFeedback}\land\text{HighControlPenalty}
\Rightarrow
P_{TMI}\le P_{std}
}
\]
То есть автоинтерфейс может ошибаться, если канал обратной связи систематически вводит его в заблуждение.

\subsection{Safety-aware следствие}

Test 11 показывает, что защитный слой уменьшает область отказа, но не устраняет её полностью:
\[
\boxed{
\text{AutoInterface}+\text{SafetyLayer}
\Rightarrow
\text{SaferAdaptiveControl}
}
\]
но не:
\[
\text{AutoInterface}+\text{SafetyLayer}
\Rightarrow
\text{PerfectControl}.
\]

\section{Что делать дальше}

Следующие шаги:
\begin{enumerate}[label=\arabic*.]
\item Подготовить независимую реализацию Test 7--8 в Qiskit/Cirq-стиле.
\item Повторить ключевые тесты на стандартном noise simulator.
\item Отделить теоретический модуль TMI-Physics 2.2 от вычислительного Supplement.
\item Подготовить репозиторий с воспроизводимой структурой: code, data, figures, reports.
\item После проверки на независимом симуляторе переходить к cloud-QPU протоколу.
\end{enumerate}

\section{Заключение}

Tests 1--8 показывают, что адаптивное интерфейсное управление может улучшать $T_2^{eff}$ и $P_{success}$ в численных моделях. Tests 9--11 делают линию научно честнее: они показывают не только область успеха, но и область слабости, отказа и частичного восстановления через safety-aware слой.

Итоговая формула всего отчёта:
\[
\boxed{
\text{Adaptive Interface Control}
\Rightarrow
\text{Measurable Computational Advantage}
}
\]
при условии корректной наблюдаемости среды, стабильного feedback-канала и safety-aware ограничений управления.

\appendix
\section{Состав пакета}

Архив отчёта содержит:
\begin{itemize}
\item PDF и LaTeX исходник данного отчёта;
\item figures/ -- ключевые графики Tests 1--11;
\item data/ -- CSV-результаты и summary-файлы;
\item test\_archives/ -- архивы отдельных Tests 1--11.
\end{itemize}

\end{document}
